% Part: first-order-logic
% Chapter: natural-deduction
% Section: provability-consistency

\documentclass[../../../include/open-logic-section]{subfiles}

\begin{document}

\iftag{FOL}
      {\olfileid[fr]{fol}{ntd}{prv}}
      {\olfileid[fr]{pl}{ntd}{prv}}
      
\olsection{\usetoken{S}{derivability} et cohérence}

Établissons plusieurs propriétés de la relation de
!!{derivability}. Outre leur intérêt propre, elles interviendront
toutes dans la démonstration du théorème de complétude.

\begin{prop}\ollabel{prop:provability-contr}
  Si $\Gamma \Proves !A$ et si $\Gamma \cup \{!A\}$ est
  incohérent, alors $\Gamma$ est incohérent.
\end{prop}

\begin{proof}
Soient $\delta_1$ une !!{derivation} de~$!A$ à partir de~$\Gamma$
et $\delta_2$ une !!{derivation} de~$\lfalse$ à partir de
$\Gamma \cup \{!A\}$. Nous obtenons alors la !!{derivation}
suivante :
\begin{prooftree}
\AxiomC{$\Gamma, \Discharge{!A}{1}$}
\RightLabel{$\delta_2$}
\DeduceC{$\lfalse$}
\DischargeRule{\Intro{\lnot}}{1}
\UnaryInfC{$\lnot !A$}
\AxiomC{$\Gamma$}
\RightLabel{$\delta_1$}
\DeduceC{$!A$}
\RightLabel{\Elim{\lnot}}
\BinaryInfC{$\lfalse$}
\end{prooftree}
Dans cette nouvelle !!{derivation}, les occurrences de
l'hypothèse~$!A$ indiquées sont !!{discharged}s. Toutes les
hypothèses qui restent non déchargées appartiennent donc à~$\Gamma$.
\end{proof}

\begin{prop}
\ollabel{prop:prov-incons}
$\Gamma \Proves !A$ si et seulement si
$\Gamma \cup \{\lnot !A\}$ est incohérent.
\end{prop}

\begin{proof}
Supposons d'abord $\Gamma \Proves !A$. Il existe donc une
!!{derivation}~$\delta_0$ de~$!A$ dont les hypothèses
!!{undischarged}s appartiennent à~$\Gamma$. Nous obtenons une
!!{derivation} de $\lfalse$ à partir de
$\Gamma \cup \{\lnot !A\}$ comme suit :
\begin{prooftree}
  \AxiomC{$\lnot !A$}
  \AxiomC{$\Gamma$}
  \RightLabel{$\delta_0$}
  \DeduceC{$!A$}
  \RightLabel{\Elim{\lnot}}
  \BinaryInfC{$\lfalse$}
\end{prooftree}
Supposons maintenant que $\Gamma \cup \{\lnot !A\}$ soit
incohérent. Soit $\delta_1$ une !!{derivation} de~$\lfalse$ dont
les hypothèses !!{undischarged}s appartiennent à
$\Gamma \cup \{\lnot !A\}$. La règle~$\FalseCl$ donne une
!!{derivation} de~$!A$ à partir de~$\Gamma$ seulement :
\begin{prooftree}
  \AxiomC{$\Gamma, \Discharge{\lnot !A}{1}$}
  \RightLabel{$\delta_1$}
  \DeduceC{$\lfalse$}
  \RightLabel{\FalseCl}
  \DischargeRule{\FalseCl}{1}
  \UnaryInfC{$!A$}
\end{prooftree}
\end{proof}

\begin{prob}
Démontrer que $\Gamma \Proves \lnot !A$ si et seulement si
$\Gamma \cup \{!A\}$ est incohérent.
\end{prob}

\begin{prop}\ollabel{prop:explicit-inc}
  Si $\Gamma \Proves !A$ et $\lnot !A \in \Gamma$, alors
  $\Gamma$ est incohérent.
\end{prop}

\begin{proof}
  Supposons $\Gamma \Proves !A$ et $\lnot !A \in \Gamma$.
  Il existe une !!{derivation}~$\delta$ de~$!A$ à partir de~$\Gamma$.
  Considérons cette application immédiate de $\Elim{\lnot}$ :
\begin{prooftree}
    \AxiomC{$\lnot !A$}
    \AxiomC{$\Gamma$}
    \RightLabel{$\delta$}
    \DeduceC{$!A$}
    \RightLabel{\Elim{\lnot}}
    \BinaryInfC{$\lfalse$}
  \end{prooftree}
Puisque $\lnot !A \in \Gamma$, toutes les hypothèses
  !!{undischarged}s appartiennent à~$\Gamma$. Cela établit
  $\Gamma \Proves \lfalse$.
\end{proof}

\begin{prop}\ollabel{prop:provability-exhaustive}
  Si $\Gamma \cup \{!A\}$ et $\Gamma \cup \{\lnot !A\}$
  sont tous deux incohérents, alors $\Gamma$ est incohérent.
\end{prop}

\begin{proof}
Il existe des !!{derivation}s $\delta_1$ et $\delta_2$ de~$\lfalse$
à partir de $\Gamma \cup \{ !A \}$ et de $\lfalse$ à partir de
$\Gamma \cup \{ \lnot !A \}$, respectivement. Nous obtenons alors :
\begin{prooftree}
\AxiomC{$\Gamma, \Discharge{\lnot !A}{2}$}
\RightLabel{$\delta_2$}
\DeduceC{$\lfalse$}
\DischargeRule{\Intro{\lnot}}{2}
\UnaryInfC{$\lnot \lnot !A$}
\AxiomC{$\Gamma, \Discharge{!A}{1}$}
\RightLabel{$\delta_1$}
\DeduceC{$\lfalse$}
\DischargeRule{\Intro{\lnot}}{1}
\UnaryInfC{$\lnot !A$}
\RightLabel{\Elim{\lnot}}
\BinaryInfC{$\lfalse$}
\end{prooftree}
Les hypothèses $!A$ et $\lnot !A$ indiquées étant
!!{discharged}s, les hypothèses encore non déchargées appartiennent
à~$\Gamma$. Nous avons donc une !!{derivation} de~$\lfalse$ à
partir de~$\Gamma$ seulement. Ainsi, $\Gamma$ est incohérent.
\end{proof}

\end{document}
